\documentclass[noanswers, addpoints, 12pt]{exam}
\usepackage[T1]{fontenc}
\usepackage[utf8]{inputenc}
\usepackage{lmodern}
\usepackage{tikz}
\usetikzlibrary{positioning,shapes.multipart,arrows.meta,shadows}
\usepackage[french]{babel}
\usepackage{xcolor}
\usepackage{listings}
\usepackage{graphicx}
\usepackage[most]{tcolorbox}

\definecolor{codegreen}{rgb}{0,0.6,0}
\definecolor{codegray}{rgb}{0.5,0.5,0.5}
\definecolor{codepurple}{rgb}{0.58,0,0.82}
\definecolor{backcolour}{rgb}{0.95,0.95,0.92}
 
\lstdefinestyle{mystyle}{
    backgroundcolor=\color{backcolour},   
    commentstyle=\color{codegreen},
    keywordstyle=\color{magenta},
    numberstyle=\tiny\color{codegray},
    stringstyle=\color{codepurple},
    basicstyle=\footnotesize,
    breakatwhitespace=false,         
    breaklines=true,                 
    captionpos=b,                    
    keepspaces=true,                 
    numbers=left,                    
    numbersep=5pt,                  
    showspaces=false,                
    showstringspaces=false,
    showtabs=false,                  
    tabsize=2
}
\lstset{style=mystyle, language=Python}

\title{CC1 1er NSI - Binaire, Variables et Conditions}
\author{Babash - NSI Tani Malandi et Johanne Bakalara}
\date{Septembre 2023, Durée: 1h}

\begin{document}
\runningfooter{}{\footnotesize Babash - NSI Tani Malandi et Johanne Bakalara \textperiodcentered{} CC BY-SA 4.0 \textperiodcentered{} page \thepage}{}
\firstpagefooter{}{\footnotesize Babash - NSI Tani Malandi et Johanne Bakalara \textperiodcentered{} CC BY-SA 4.0 \textperiodcentered{} page \thepage}{}

\maketitle

\pagestyle{headandfoot}
\thispagestyle{headandfoot}
\firstpageheader{Nom : \\  Prénom :\\ Classe :}{}{}


\section*{Consignes:}
\begin{center} 
\fbox{\parbox{\linewidth}{
\begin{center}
Le devoir est à rédiger sur une feuille libre.\\
La clarté de la rédaction et les explications font partie intégrante de la notation. Une réponse non rédigée ou sans explication sera considérée fausse.\\
Prenez le temps d'encadrer les réponses.
\end{center}


}}
\end{center}


\vsword{Note}
\hsword{Note}

\section*{Barème :}
\begin{center}
  \par\bigskip
  \gradetable[h]
\end{center}


\section{Conversion Octet-bits: }

\begin{questions}
\question[1] Sam dispose d'un téléphone contenant 80 Go de mémoire, il aime bien écouter de la musique durant les pauses.
Une musique pèse environ 240 Mbits. Combien Sam peut-il stocker de musique sur son téléphone ?
\begin{solution}[5\bigskipamount]
     Commençons par convertir les 80 Go en bits, on sait que  $1\ Go = 10^9\ o$ donc $80\ Go = 80 \times 10^9\ o$ par ailleurs, on sait que $1\ octet = 8\ bits $ donc $80\ Go = 80 \times 10^9\ o = 80 \times 10^9 \times 8\ bits$.\\

     Convertissons maintenant les $240\ Mbits$, on sait que $1\ Mbits = 10^6\ bits$ donc $240\ Mbits = 240 \times 10^6\ bits$.\\

     Maintenant, on souhaite connaitre le nombre de musiques qui peuvent rentrer sur le téléphone, pour cela, on doit effectuer une division, on a donc\ :
     $$\frac{80 \times 10^9 \times 8\ bits}{240 \times 10^6\ bits} = 10^3 \times \frac{80 \times 8}{240} = 10^3 \times \frac{8}{3} \approx 2666{,}66$$  On déduit qu'il y a de la place pour 2666 musiques complètes sur le téléphone.
\end{solution}

\question[1] Combien de valeurs différentes peut-on stocker sur 16 bits.
\begin{solution}[5\bigskipamount]
     Sur 16 bits, on peut écrire les nombres de $0$ à $2^{16} - 1$ on peut donc stocker $2^{16} = 65536$ valeurs différentes.
\end{solution}

\section{Décimal, binaire et hexadécimal}

\begin{tcolorbox}[enhanced,
    colback=yellow!25!black!10!white,colframe=yellow!75!black,title=Rappel,
    drop fuzzy shadow,watermark color=white]
    \begin{itemize}
        \item On appelle la base 2 la base binaire, elle est composée uniquement de 2 symboles 0 et 1.
        \item On appelle la base 10 la base décimale, elle est composée de 10 symboles qui sont les chiffres de 0 à 9.
        \item On appelle la base 16 la base hexadécimale, elle est composée de 16 symboles qui sont les chiffres de 0 à 9 puis les lettres de A à F.
    \end{itemize}
\end{tcolorbox}

\question[1] Convertissez $63_{10}$ en binaire (vous détaillerez vos calculs). Expliquez pourquoi ce résultat était facilement calculable.
\begin{solution}[5\bigskipamount]
     Pour passer de la base 10 à la base 2 on effectue une décomposition du nombre en puissance de 2.
     
     Ici la plus grande puissance de 2 inférieure à 63 est $2^5 = 32$, on peut donc écrire $63_{(10)} = 2^5 + 31$.
     
     Il nous reste donc à décomposer $31$, la puissance de 2 suivante est $2^4=16$ qui est inférieure à 31, on écrit donc $63_{(10)} = 2^5 + 2^4 + 15$.
     
     En procédant de la même manière avec les puissances suivantes, on a $63_{(10)} = 2^5 + 2^4 + 2^3 + 2^2 + 2^1 + 2^0$

     En reportant sur notre tableau de puissance de 2 :
    \begin{center}
     \begin{tabular}{lcccccccc}
        & $2^7$ & $2^6$ & $2^5$ & $2^4$ & $2^3$ & $2^2$ & $2^1$ & $2^0$\\
        & $128$ & $64$ & $32$ & $16$ & $8$ & $4$ & $2$ & $1$\\
        \hline
        $63_{(10)} = $ & 0 & 0 & 1 & 1 & 1 & 1 & 1 & 1\\
    \end{tabular}
    \end{center}

    On a donc $63_{(10)} = 11\ 1111_{(2)}$

    Ce résultat était prévisible, car $63_{(10)} = 64 - 1 = 2^6 - 1 = 100\ 0000_{(2)} - 1 = 11\ 1111_{(2)}$
\end{solution}

\question[1] Convertir le nombre $113_{16}$ en décimal.
\begin{solution}[5\bigskipamount]
     Pour passer de la base 16 à la base 10 on peut utiliser le tableau de la base 16 :
     \begin{center}
     \begin{tabular}{lccc}
        & $16^2$ & $16^1$ & $16^0$\\
        & $256$ & $16$ & $1$\\
        \hline
        $113_{(16)} = $ & 1 & 1 & 3\\
    \end{tabular}
    \end{center}

    On a donc la décomposition : $113_{(16)} = 16^2 + 16^1 + 3 = 256 + 16 + 3 = 275_{(10)}$ 
\end{solution}

\begin{tcolorbox}[enhanced,
    colback=yellow!25!black!10!white,colframe=yellow!75!black,title=Rappel,
    drop fuzzy shadow,watermark color=white]
    % Table ASCII (codes 32 à 126), dessinée pour ce document.
\begin{center}
\footnotesize
\setlength{\tabcolsep}{3pt}
\renewcommand{\arraystretch}{1.3}
\begin{tabular}{|c|*{16}{c|}}
\hline
\textbf{hex} & \textbf{0} & \textbf{1} & \textbf{2} & \textbf{3} & \textbf{4} & \textbf{5} & \textbf{6} & \textbf{7} & \textbf{8} & \textbf{9} & \textbf{A} & \textbf{B} & \textbf{C} & \textbf{D} & \textbf{E} & \textbf{F} \\ \hline
\textbf{2} & \textsc{esp} & ! & " & \# & \$ & \% & \& & ' & ( & ) & * & + & , & - & . & / \\ \hline
\textbf{3} & 0 & 1 & 2 & 3 & 4 & 5 & 6 & 7 & 8 & 9 & : & ; & < & = & > & ? \\ \hline
\textbf{4} & @ & A & B & C & D & E & F & G & H & I & J & K & L & M & N & O \\ \hline
\textbf{5} & P & Q & R & S & T & U & V & W & X & Y & Z & [ & \textbackslash & ] & \^{} & \_ \\ \hline
\textbf{6} & ` & a & b & c & d & e & f & g & h & i & j & k & l & m & n & o \\ \hline
\textbf{7} & p & q & r & s & t & u & v & w & x & y & z & \{ & | & \} & \textasciitilde & \textsc{del} \\ \hline
\end{tabular}

\smallskip
Lecture : le caractère \texttt{A} est à la ligne \textbf{4}, colonne \textbf{1} : son code ASCII est \texttt{41} en hexadécimal ($65$ en décimal).
\end{center}

\end{tcolorbox}
\question[2] Décoder le message suivant en ASCII (Le plus simple est de passer par la base 16) : 
01010110 01101111 01110101 01110011 00100000 01101101 01100001 01110010 01110001 01110101 01100101 01111010 00100000 00110010 00100000 01110000 01110100 01110011 00100000 00100001

\begin{solution}[5\bigskipamount]
    Pour décoder le message, on doit d'abord traiter chaque symbole indépendamment. En ASCII, les symboles sont représentés sur 8 bits (1 octet) chaque nombre entre 0 et 127 se voit donc associer un symbole.
    Pour connaître ce symbole, on utilise le tableau, mais il nous faut d'abord convertir le nombre soit en base 10 soit en base 16 comme la conversion de la base 2 dans la base 16 est très facile, on choisit cette technique.

    Pour le premier symbole, le code binaire est $01010110$ on regroupe alors les bits par 4, on a donc $0101\ 0110$.
    
    On convertit ensuite chaque bloc séparément :
    \begin{itemize}
        \item $0101 = 2^2 + 2^0 = 4 + 1 = 5$
        \item $0110 = 2^2 + 2^1 = 4 + 2 = 6$
    \end{itemize}
    On déduit $01010110_{(2)} = 56_{(16)}$, d'après le tableau, le premier symbole est donc un V majuscule.

    En procédant de la même manière pour tous les symboles, on obtient le message : "Vous marquez 2 pts !".

    Mais ne vous inquiétez pas, si vous faites tout ça vous en marquez bien 2 =p
\end{solution}


\section{Complément à 2}

\question[3] Question de cours et réflexion
\begin{parts}
\part À quoi sert le complément à 2 ?
\begin{solution}[5\bigskipamount]
    Le complément à 2 est une méthode permettant de stocker les entiers relatifs (et donc en particulier les négatifs) sur un ordinateur.
    
    Elle a l'avantage de conserver l'algorithme d'addition classique et la même représentation pour les nombres positifs.    
\end{solution}

\part Parmi ces 4 propositions (phrases) dites lesquelles sont vraies et lesquelles sont fausses :
\begin{itemize}
    \item En complément à 2 les nombres positifs gardent leurs représentations classiques.
    \item En complément à 2 l'addition d'un nombre et de son opposé vaut 0.
    \item En complément à 2 le zéro a deux représentations, une positive, une négative.
    \item En complément à 2 l'algorithme d'addition reste le même.
\end{itemize}
\begin{solution}[5\bigskipamount]
    \begin{itemize}
        \item Vrai
        \item Vrai
        \item Faux
        \item Vrai
    \end{itemize}
\end{solution}

\part Par quels symboles commencent les nombres positifs et les nombres négatifs en compléments à 2.
\begin{solution}[5\bigskipamount]
    Les nombres positifs commencent par un 0, les nombres négatifs par un 1.    
\end{solution}
\end{parts}

\question[1] Quelle est la valeur décimale du nombre $1001 0010_{(2)}$ ici représenté en complément à 2 ?
\begin{solution}[5\bigskipamount]
    On peut utiliser le tableau du complément à 2 :
    \begin{center}
     \begin{tabular}{lcccccccc}
        & $-2^7$ & $2^6$ & $2^5$ & $2^4$ & $2^3$ & $2^2$ & $2^1$ & $2^0$\\
        & $-128$ & $64$ & $32$ & $16$ & $8$ & $4$ & $2$ & $1$\\
        \hline
        & 1 & 0 & 0 & 1 & 0 & 0 & 1 & 0\\
    \end{tabular}
    \end{center}

    On a donc la décomposition : $1001 0010_{(2)} = -2^7 + 2^4 + 2 = -128 + 16 + 2 = -110$. 
\end{solution}

\section{Variables}

\section{Conditions}

\end{questions}
\end{document}

